\documentclass[a4paper]{article}

%------------------------------------------------
% الحزم
%------------------------------------------------

\usepackage{enumitem}
\usepackage[margin=2.5cm]{geometry}
\usepackage{amsmath,amsfonts,amssymb}
\usepackage[charter]{mathdesign}
\usepackage{graphicx}
\usepackage{array}
\usepackage{fancyhdr}
\usepackage{polyglossia}

%------------------------------------------------
% اللغات والخطوط
%------------------------------------------------

\setdefaultlanguage[
    calendar=gregorian,
    numerals=maghrib
]{arabic}

\setotherlanguage{english}

% إذا لم تكن الخطوط مثبتة في الجهاز، تُؤخذ من المجلد ../fonts/
\IfFontExistsTF{Amiri}{%
    \newfontfamily\arabicfont[Script=Arabic, Scale=1.15]{Amiri}
}{%
    \newfontfamily\arabicfont[
        Script=Arabic, Scale=1.15,
        Path=../fonts/, Extension=.ttf,
        UprightFont=*-Regular, BoldFont=*-Bold,
        ItalicFont=*-Italic, BoldItalicFont=*-BoldItalic
    ]{Amiri}
}

\IfFontExistsTF{AlBattar}{%
    \newfontfamily\arabicfontsf[Script=Arabic, Scale=1.2]{AlBattar}
    % خط العناوين (يعمل حتى بعد استعمال \textenglish)
    \newfontfamily\khatt[Script=Arabic, Scale=1.2]{AlBattar}
}{%
    \newfontfamily\arabicfontsf[
        Script=Arabic, Scale=1.2, Path=../fonts/,
        BoldFont=ae_AlBattar.ttf
    ]{ae_AlBattar.ttf}
    % خط العناوين (يعمل حتى بعد استعمال \textenglish)
    \newfontfamily\khatt[
        Script=Arabic, Scale=1.2, Path=../fonts/,
        BoldFont=ae_AlBattar.ttf
    ]{ae_AlBattar.ttf}
}

%------------------------------------------------
% إعداد الصفحة
%------------------------------------------------

\renewcommand{\footrulewidth}{1pt}
\renewcommand{\headrulewidth}{1pt}

\rhead{\khatt المدرسة العليا للأساتذة ببوزريعة}
\lhead{\khatt ملحقة بومرداس}

\cfoot{}
\lfoot{\khatt الأستاذ: د. بن عبد الرزاق يوسف}

\parindent 0pt
\setlength{\headsep}{15pt}
\setlength{\parskip}{3pt}

\pagestyle{fancy}

\setlist[enumerate,1]{
    label=\arabic*),
    itemsep=5pt,
    topsep=3pt
}

\setlist[itemize]{
    itemsep=3pt,
    topsep=3pt
}

\arraycolsep=3pt

%================================================
% بداية الوثيقة
%================================================

\begin{document}



\rule{\textwidth}{1pt}

\begin{center}

{\khatt\textbf{
الفصل 6 : الذكاء الاصطناعي متعدد الأنماط
}}

\end{center}

\rule{\textwidth}{1pt}





%================================================
% النشاط الأول
%================================================

\section*{\khatt\underline{النشاط الأول: فهم صورة باستعمال الذكاء الاصطناعي}}

\begin{enumerate}

%------------------------------------------------

\item اختر صورة بسيطة تحتوي على معلومات، مثل:

\begin{itemize}

    \item مخطط بياني.

    \item جدول.

    \item صورة علمية.

    \item رسم توضيحي.

\end{itemize}

ثم ارفع الصورة إلى إحدى أدوات الذكاء الاصطناعي
متعددة الأنماط، مثل:

\textenglish{ChatGPT} أو \textenglish{Gemini} أو
\textenglish{Claude}.

%------------------------------------------------

\item اكتب أمرًا يطلب من الأداة فهم الصورة وشرحها.

مثال:

\begin{center}

\textbf{
حلل هذه الصورة واشرح لي المعلومات الموجودة فيها
بطريقة بسيطة
}

\end{center}

اكتب الأمر الذي استعملته:



%------------------------------------------------

\item اكتب باختصار ماذا فهمت الأداة من الصورة:

%------------------------------------------------

\item قارن بين إجابة الذكاء الاصطناعي والصورة الأصلية:

\begin{center}

\begin{tabular}{|
>{\centering\arraybackslash}m{4cm}|
>{\centering\arraybackslash}m{5cm}|
>{\centering\arraybackslash}m{5cm}|}

\hline

\textbf{العنصر}
&
\textbf{إجابة الذكاء الاصطناعي}
&
\textbf{المعلومة الصحيحة في الصورة}
\\

\hline

الأرقام
&
&
\\[0.7cm]

\hline

المعلومات
&
&
\\[0.7cm]

\hline

العلاقات بين العناصر
&
&
\\[0.7cm]

\hline

\end{tabular}

\end{center}

%------------------------------------------------

\item هل كانت إجابة الذكاء الاصطناعي صحيحة؟

ضع علامة:

\begin{center}

$\square$ نعم

\hspace{2cm}

$\square$ لا

\hspace{2cm}

$\square$ كانت صحيحة جزئيًا

\end{center}



\end{enumerate}


%================================================
% النشاط الثاني
%================================================

\section*{\khatt\underline{النشاط الثاني (واجب منزلي): تحويل النص إلى مخطط}}

\textbf{يُنجز هذا النشاط في المنزل ويُسلَّم في الحصة القادمة.}

\begin{enumerate}

%------------------------------------------------

\item اختر فقرة قصيرة من إحدى موادك الدراسية.

يفضل أن تحتوي الفقرة على عدة مفاهيم مرتبطة ببعضها.

%------------------------------------------------

\item أعط النص لأداة ذكاء اصطناعي واطلب منها
تحويله إلى مخطط مفاهيمي.

مثال:

\begin{center}

\textbf{
``حوّل هذا النص إلى مخطط مفاهيمي بسيط يوضح
العلاقة بين المفاهيم الرئيسية.''
}

\end{center}

%------------------------------------------------

\item اكتب الأمر الذي استعملته:


%------------------------------------------------

\item ارسم المخطط الذي اقترحته الأداة في المساحة التالية:

\begin{center}

\fbox{
\parbox[c][5cm][c]{13cm}{
\centering
\textbf{ }
}}

\end{center}

%------------------------------------------------

\item هل يحتوي المخطط على أهم أفكار النص؟

\begin{center}

$\square$ نعم

\hspace{2cm}

$\square$ لا

\hspace{2cm}

$\square$ بعض الأفكار فقط

\end{center}

%------------------------------------------------

\item هل توجد علاقة غير صحيحة بين بعض المفاهيم؟

إذا نعم، اذكرها:



%------------------------------------------------

\item هل أضفت أو حذفت شيئًا من المخطط؟

اشرح باختصار:



\end{enumerate}


\end{document}